\documentclass[11pt]{article}

\usepackage{geometry}
\usepackage{amsmath}
\usepackage{enumitem}
\usepackage{hyperref}

\geometry{margin=1in}

\title{CNV Segment to Gene Set Methods}
\author{}
\date{}

\begin{document}
\maketitle

\section{Overview}

The CNV extractor (\texttt{cnv\_gene\_extractor}) maps CNV segments to genes,
computes signed gene scores, and emits directional gene programs:
amplification (\texttt{amp}) and deletion (\texttt{del}).

The implementation is designed for segment-level inputs (tumor or germline)
without requiring integer copy-number reconstruction.

\section{Inputs and notation}

For each sample $s$, let segments be indexed by $i$ with:
\[
(\mathrm{chrom}_i, \mathrm{start}_i, \mathrm{end}_i, a_i),
\]
where $a_i$ is the signed segment amplitude (for example log2 copy-ratio).

Common segment-header presets are supported:
\texttt{segments\_format=gdc\_seg} and \texttt{segments\_format=cbio\_seg}
(cBioPortal SEG: \texttt{ID, chrom, loc.start, loc.end, seg.mean}).

Genes from the GTF are indexed by $g$ with gene-body interval
$[\mathrm{gstart}_g, \mathrm{gend}_g)$ on chromosome $\mathrm{chrom}_g$.

\section{Segment-to-gene overlap}

For segment $i$ and gene $g$, overlap length is:
\[
\mathrm{ovbp}_{ig}
=
\max\left(0,\,
\min(\mathrm{end}_i,\mathrm{gend}_g) -
\max(\mathrm{start}_i,\mathrm{gstart}_g)\right).
\]

Gene-length-normalized overlap fraction:
\[
f_{ig} = \frac{\mathrm{ovbp}_{ig}}{\max(1,\mathrm{gend}_g-\mathrm{gstart}_g)}.
\]

\section{Focality and breadth penalties}

Segment length:
\[
L_i = \max(1, \mathrm{end}_i - \mathrm{start}_i).
\]

Focal penalty:
\[
p^{\text{focal}}_i
=
\frac{1}{1 + \left(\frac{L_i}{L_{\text{scale}}}\right)^{\alpha}},
\]
with CLI controls \texttt{--focal\_length\_scale\_bp} and
\texttt{--focal\_length\_alpha}.

Gene-count penalty for $n_i$ genes overlapping segment $i$:
\[
p^{\text{genes}}_i =
\begin{cases}
1 & \text{mode=none}\\
\frac{1}{\sqrt{n_i}} & \text{mode=inv\_sqrt}\\
\frac{1}{n_i} & \text{mode=inv}
\end{cases}
\]
controlled by \texttt{--gene\_count\_penalty}.

Combined segment weight:
\[
w_i = p^{\text{focal}}_i \cdot p^{\text{genes}}_i.
\]

Preset interpretation:
\begin{itemize}[leftmargin=2em]
\item \texttt{default}: balanced focal penalties.
\item \texttt{focal}: stronger focal penalties (and optional segment-length cap) to emphasize focal events.
\item \texttt{broad}: weaker penalties for arm-level/aneuploidy-oriented summaries.
\item \texttt{qc}: emits both focal and broad variants per sample for side-by-side interpretation.
\end{itemize}

\section{Gene scoring}

Default per-gene numerator:
\[
N_g = \sum_i a_i \cdot w_i \cdot f_{ig}.
\]

Default aggregation (\texttt{weighted\_mean}):
\[
\mathrm{score}_g = \frac{N_g}{\sum_i |f_{ig}|}.
\]

Alternative aggregations (\texttt{sum}, \texttt{mean}, \texttt{max\_abs})
are exposed by \texttt{--aggregation}.

Segments with $|a_i| < \texttt{--min\_abs\_amplitude}$ are dropped before scoring.

\section{Directional program construction}

From signed gene score $\mathrm{score}_g$:
\[
\mathrm{amp}_g = \max(\mathrm{score}_g, 0),
\qquad
\mathrm{del}_g = \max(-\mathrm{score}_g, 0).
\]

These define \texttt{program=amp} and \texttt{program=del}.
Optional \texttt{program=abs} uses $|\mathrm{score}_g|$.

\section{Selection, normalization, and GMT}

Per program, the extractor applies the shared selection operator
(\texttt{top\_k}, \texttt{quantile}, or \texttt{threshold}) and then
normalizes selected weights (default \texttt{within\_set\_l1}).

Default GMT behavior targets compact, enrichment-ready sets:
\begin{itemize}[leftmargin=2em]
\item \texttt{--gmt\_topk\_list 200}
\item \texttt{--gmt\_min\_genes 100}
\item \texttt{--gmt\_max\_genes 500}
\item sets below minimum are skipped unless \texttt{--emit\_small\_gene\_sets true}
\end{itemize}

For \texttt{top\_k}, ordering is deterministic (score, then gene symbol, then gene ID).
If the cutoff intersects a tie block, tie diagnostics are recorded in metadata and
the flat program manifest.

\section{Optional purity correction}

If enabled (\texttt{--use\_purity\_correction true}), amplitude is adjusted as:
\[
a'_i = \frac{a_i}{\max(\mathrm{purity}_s, p_{\min})},
\]
with caps via \texttt{--max\_abs\_amplitude}.

This is a lightweight correction and not a full ploidy model.

\section{Optional cohort recurrence sets}

Given sample-level scores, recurrent amp/del genes are defined by fraction
of samples passing a score threshold:
\[
r^{\text{amp}}_g = \frac{1}{|S|}\sum_{s \in S} \mathbf{1}\left[\max(\mathrm{score}_{g,s},0) > \tau\right],
\]
\[
r^{\text{del}}_g = \frac{1}{|S|}\sum_{s \in S} \mathbf{1}\left[\max(-\mathrm{score}_{g,s},0) > \tau\right].
\]
Genes passing \texttt{--cohort\_min\_fraction} are ranked by recurrence.

\section{Practical caveats}

\begin{itemize}[leftmargin=2em]
\item Broad CNV burden can dominate interpretation; defaults penalize broad segments.
\item Cytoband and region-level enrichment can dominate CNV enrichment outputs and is expected as positional QC; evaluate mechanism-oriented interpretations separately.
\item If arm-level aneuploidy is expected, use \texttt{--program\_preset broad}.
\item If focal driver-like events are expected, use \texttt{--program\_preset focal}.
\item Coordinate-system and chromosome-prefix mismatches can remove many segments.
The implementation runs a chromosome-compatibility preflight check before scoring and warns (or fails) when compatibility is very low.
\item Purity correction can amplify noise in low-purity samples; review run summaries.
\item If a requested program has no positive signal after directional projection, it is skipped and recorded in \texttt{skipped\_programs.json}; per-program QC rows are written to \texttt{cnv\_program\_manifest.tsv}. Manifest rows explicitly separate \texttt{scored\_ok} from \texttt{gmt\_emitted} to avoid treating score-only rows as emitted GMT outputs.
\end{itemize}

\end{document}
